\documentclass[../../main.tex]{subfiles}
\begin{document}

{\bf[解]}

\begin{figure}[htb]
  \centering
  % Idisso ti panagkita iti 3D (tdplotsetmaincoords{theta}{phi})
  \tdplotsetmaincoords{70}{115}
  \begin{tikzpicture}[tdplot_main_coords, scale=2.4, >=stealth, every node/.style={font=\small}]

    % --- 1. Koordinada dagiti uppat a vertex (Tetrahedron) ---
    % A0 ti tuktok (apex), A1, A2, A3 ti base
    \coordinate (A0) at (0, 0, 1.6);
    \coordinate (A1) at ({1.3*cos(210)}, {1.3*sin(210)}, 0);
    \coordinate (A2) at ({1.3*cos(330)}, {1.3*sin(330)}, 0);
    \coordinate (A3) at ({1.3*cos(90)},  {1.3*sin(90)},  0);

    % --- 2. Kolor ti rabaw (Faces shading) ---
    % Rabaw iti likud: A1-A2-A3 (base)
    \fill[blue!5, opacity=0.6] (A1) -- (A2) -- (A3) -- cycle;
    % Rabaw iti sanguanan: A0-A1-A2
    \fill[blue!15, opacity=0.6] (A0) -- (A1) -- (A2) -- cycle;
    % Rabaw iti kanigid: A0-A1-A3
    \fill[blue!10, opacity=0.4] (A0) -- (A1) -- (A3) -- cycle;

    % --- 3. Dagiti igid (Edges) ---
    % Nailemmeng nga igid iti likud (dashed lines)
    \draw[thick, dashed, gray!70] (A1) -- (A3);
    \draw[thick, dashed, gray!70] (A2) -- (A1);
    \draw[thick, dashed, gray!70] (A0) -- (A1);

    % Nalitnaw nga igid iti sanguanan (solid lines)
    \draw[very thick, black] (A0) -- (A2) -- (A3) -- cycle;

    % --- 4. Dagiti pana ti transision (Probability 1/3) ---
    % Manipud A0 agturong iti A1
    \draw[->, thick, red!85!black]
    ($ (A0)!0.25!(A1) $) -- ($ (A0)!0.75!(A1) $)
    node[midway, above left=1pt, font=\footnotesize] {$\dfrac{1}{3}$};

    % Manipud A0 agturong iti A2
    \draw[->, thick, red!85!black]
    ($ (A0)!0.25!(A2) $) -- ($ (A0)!0.75!(A2) $)
    node[midway, above right=1pt, font=\footnotesize] {$\dfrac{1}{3}$};

    % Manipud A0 agturong iti A3 (dashed ta adda iti likud)
    \draw[->, thick, dashed, red!70!black]
    ($ (A0)!0.25!(A3) $) -- ($ (A0)!0.75!(A3) $)
    node[midway, right=2pt, font=\footnotesize] {$\dfrac{1}{3}$};

    % --- 5. Dagiti marka ken nagan ti Vertex ---
    % A0 (Irugi iti t = 0)
    \fill[red!85!black] (A0) circle (1.0pt);
    \node[above=3pt, red!85!black, font=\bfseries] at (A0) {$A_0$ ($t=0$)};

    % A1, A2, A3
    \fill (A1) circle (0.8pt) node[below left=2pt] {$A_1$};
    \fill (A2) circle (0.8pt) node[below right=2pt] {$A_2$};
    \fill (A3) circle (0.8pt) node[above=2pt] {$A_3$};

  \end{tikzpicture}
  \caption{四面体 $A_i$ 上の各頂点間の遷移}
  \label{fig:tetrahedron_random_walk}
\end{figure}

まず$p_i(n+1)$について考えると，$t=n+1$で頂点$i$にいるのは，$t=n$で頂点$i$以外にいて，そこから確率$\dfrac{1}{3}$で頂点$i$に来た時だから，
\begin{align}
  p_i(n+1)=\dfrac{1}{3}(1-p_i(n)) \label{eq:1}
\end{align}
となる．
初期条件は題意より，$t=0$において頂点$i=0$にあることから
\begin{align}
  p_0(0) = 1 \quad p_{k}(0)=0\, (k=1,2,3) \label{eq:2}
\end{align}
である．

\subparagraph{(1)}

数学的帰納法で$p_1(n)=p_2(n)=p_3(n)$を証明する．
\cref{eq:2}から$n=0$では題意が成立する．
そこで$n=k$での成立を仮定すれば，\cref{eq:1}より自明に$n=k+1$でもこの条件は成立する．
従って数学的帰納法により．
\begin{align}
  p_1(n)=p_2(n)=p_3(n)
\end{align}
である．

\subparagraph{(2)}

\cref{eq:1}を変形して
\begin{align}
  p_i(n+1)-\dfrac{1}{4}=\dfrac{-1}{3}\left(p_i(n)-\dfrac{1}{4}\right) \label{eq:3}
\end{align}
だから繰り返し用いて
\begin{align}
  p_i(n)=\left(\dfrac{-1}{3}\right)^n\left(p_i(0)-\dfrac{1}{4}\right)+\dfrac{1}{4}
\end{align}
である．
初期条件は\cref{eq:2}より$p_1(0)=0$，$p_0(1)=1$だから代入して
\begin{align}
  \begin{cases}
    p_0(n)=\dfrac{3}{4}\left(\dfrac{-1}{3}\right)^n+\dfrac{1}{4} \\
    p_1(n)=\dfrac{1}{4}\left\{1-\left(\dfrac{-1}{3}\right)^n\right\}
  \end{cases}
\end{align}
が求める値である．$\cdots$(答)
\newpage
\end{document}
